\scanpage{32}
\unsure{16|15}) \uline{Classes fondamentales globales}

On fixe un entier $n>0$, et on se
\deleted{Soit $f:X\mathbin{-}Y$ morphisme}
borne à des faisceaux de $n$-torsion (i.e. Modules sur $A=\mathbf Z/n\mathbf Z$), et des préschémas à car. résiduelles premières à $n$ i.e tels que $n.1_{\mathcal O_X}$ soit inversible dans $\mathcal O_X$.

Si $f:X\mathbin{-}Y$ est un morphisme lisse, on dira qu'il est \inserted{purement} de dim relative $d$ si ses fibres sont \inserted{pures} de dim $d$. On \inserted{désigne} par $T_{X/Y}^{(n)}$ ou simplement $T_{X/Y}$ s'il n'y a pas de confusion possible, le faisceau $\mu_{nX}^{\otimes d}$. [Si $f$ est lisse de dim \inserted{relative \unsure{loc\textsuperscript{t}} variable} quelconque, alors on définit de façon évidente $T_{X/Y}$ en le prenant sur ses divers morceaux].

Nous regardons pour la suite des morphismes \uline{lisses} \uline{compactifiables} \uline{purs} \uline{de dim} $d$\quad $f:X\longrightarrow Y$.

Pour un tel morphisme, on peut définir un \uline{\deleted{\unsure{trace|hom}} \inserted{hom.} canonique}
\mspage{31}
